\documentclass[11pt,a4paper]{article}

%% PACHETE MATEMATICE
\usepackage{amsmath,amssymb,amsfonts}
\usepackage{amsthm}
\usepackage{mathpazo}
\usepackage{eulervm}
\usepackage{enumerate}
\usepackage{color}
\usepackage{graphicx}
\usepackage[all]{xy}
\usepackage{xcolor}
\usepackage{framed}
\usepackage[unicode]{hyperref}
\setcounter{MaxMatrixCols}{20}
\usepackage{multicol}
%\usepackage[romanian]{babel}


%% ASPECT
\linespread{1.05}
\usepackage[T1]{fontenc}
\usepackage[utf8x]{inputenc}
\usepackage[margin=2cm]{geometry}
% \usepackage[color]{showkeys}
\usepackage{anyfontsize}
\usepackage{titlesec}
\titleformat*{\section}{\large\bfseries}

\usepackage{fancyhdr}
\pagestyle{fancy}
\rhead{\small \color{gray} Notițe de Adrian Manea}
\lhead{\small \color{gray} IntProb-FIA, 2017---2018, L. Lemnete \& M. Olteanu}


\usepackage[autostyle = false, style = german]{csquotes}
\usepackage[romanian]{babel}
\newcommand\qq{\enquote}

%% COMENZILE MELE
\renewcommand*{\proofname}{Dem.:}
\newcommand\bb{\mathbb}
\newcommand\dr{\mathrm}
\newcommand\kal{\mathcal}
\newcommand\fk{\mathfrak}
\newcommand\op{\oplus}
\newcommand\ot{\otimes}
\newcommand\ds{\displaystyle}
\newcommand\id{\indent}
\newcommand\nid{\noindent}
\newcommand\seq{\subseteq}
\newcommand\sneq{\subsetneq}
\newcommand\speq{\supseteq}
\newcommand\spneq{\supsetneq}
\newcommand\rar{\rightarrow}
\newcommand\Rar{\Rightarrow}
\newcommand\xrar{\xrightarrow}
\newcommand\lar{\leftarrow}
\newcommand\Lar{\Leftarrow}
\newcommand\xlar{\xleftarrow}
\newcommand\lrar{\leftrightarrow}
\newcommand\Lrar{\Leftrightarrow}
\newcommand\ideal{\trianglelefteq}
\newcommand\ai{\text{ a.\^{i}. }}
\newcommand\sk{(\textit{Schi\c{t}\u{a}}) }
\newcommand\rhup{\rightharpoonup}
\newcommand\rhdn{\rightharpoondown}
\newcommand\lhup{\leftharpoonup}
\newcommand\lhdn{\leftharpoondown}
\newcommand\vid{\emptyset}
\newcommand\wh{\widehat}
\newcommand\ol{\overline}
\newcommand\NN{\mathbb{N}}
\newcommand\RR{\mathbb{R}}
\newcommand\ZZ{\mathbb{Z}}
\newcommand\QQ{\mathbb{Q}}
\newcommand\CC{\mathbb{C}}

%% STILURILE MELE
\newtheoremstyle{dotless}{}{}{\itshape}{}{\bfseries}{:}{ }{}

\theoremstyle{dotless}
\newtheorem{theorem}{Teorem\u{a}}[section]
\newtheorem{proposition}{Propozi\c{t}ie}[section]
\newtheorem{lemma}{Lem\u{a}}[section]
\newtheorem{corollary}{Corolar}[section]
\newtheorem{exercise}{Exerci\c{t}iu}

\newtheoremstyle{dotlessdr}{}{}{}{}{\bfseries}{:}{ }{}
\theoremstyle{dotlessdr}
\newtheorem{example}{Exemplu}[section]
\newtheorem{definition}{Defini\c{t}ie}[section]
\newtheorem{remark}{Observa\c{t}ie}[section]




\begin{document}

\begin{center}
  {\large\textbf{Seminar 1 \\ Integrale multiple}}
\end{center}


\vspace{2cm}


1.  Calculați:
\begin{multicols}{2}
  \begin{enumerate}[(a)]
  \item $\ds\int \frac{2x + 1}{\sqrt{x^2 - 16}} dx, x > 4$;
  \item $\ds\int \frac{2x - 5}{x^2 - 5x + 7} dx$;
  \item $\ds\int \frac{x - 1}{3x^2 - 6x + 11} dx$;
  \item $\ds\int \frac{2x}{x^4 - 1} dx, x \in (-1, 1)$;
  \item $\ds\int \arcsin x dx$;
  \item $\ds\int_1^{\sqrt{e}} x \log_3 x dx$;
  \item $\ds\int_1^e \sin(\ln x) dx$;
  \item $\ds\int_0^{\frac{1}{\sqrt{2}}} \frac{3 - 2x}{2x^2 + 1} dx$;
  \item $\ds\int_1^e \frac{\ln x}{x(2 + \ln x)} dx$;
  \end{enumerate}
\end{multicols}

Indicații:
\begin{itemize}
\item $\ds\int \frac{dx}{x^2 + a^2} = \frac{1}{a} \arctan \frac{x}{a} + c$;
\item $\ds\int \frac{dx}{x^2 - a^2} = \frac{1}{2a} \ln \Big| \frac{x - a}{x + a} \Big| + c$;
\item $\ds\int \frac{dx}{\sqrt{x^2 + a^2}} = \ln(x + \sqrt{x^2 + a^2}) + c$;
\item $\ds\int \frac{dx}{\sqrt{x^2 - a^2}} = \ln|x + \sqrt{x^2 - a^2}| + c$;
\item $\ds\int \frac{dx}{\sqrt{a^2 - x^2}} = \arcsin \frac{x}{a} + c$;
\item $\ds\int \frac{-dx}{\sqrt{a^2 - x^2}} = \arccos \frac{x}{a} + c$;
\end{itemize}


\vspace{1cm}


2. Să se calculeze integralele duble:
\begin{enumerate}[(a)]
\item $\ds\iint_D xy^2 dxdy$, unde $D = [0, 1] \times [2, 3]$;
\item $\ds\iint_D xy dxdy$, unde $D = \{(x, y) \in \RR^2 \mid y \in [0, 1], y^2 \leq x \leq y\}$;
\item $\ds\iint_D (x + 3y) dxdy$, unde $D$ este mulțimea plană mărginită de curbele
  de ecuații $y = x^2 + 1, y = -x^2, x = -1, x = 3$;
\item $\ds\iint_D xdxdy$, unde $D$ este mulțimea plană mărginită de $x^2 + y^2 = 9, x \geq 0$;
\item $\ds\iint_D xy dxdy$, unde $D$ este domeniul delimitat de curbele $xy = 1$ și
  $x + y = \frac{5}{2}$;
\item $\ds\iint_D ydxdy$, unde $D$ este domeniul mărginit de parabola $y^2 = 2x$,
  cercul $x^2 + y^2 - 2x = 0$ și dreapta $x = 2$.
\end{enumerate}

\vspace{1cm}

3. Să se calculeze, trecînd la coordonate polare, integralele:
\begin{enumerate}[(a)]
\item $\ds\iint_D e^{x^2 + y^2} dxdy$, unde $D = \{(x, y) \in \RR^2 \mid x^2 + y^2 \leq 1 \}$;
\item $\ds\iint_D 1 + \sqrt{x^2 + y^2} dxdy$, unde $D = \{(x, y) \in \RR^2 \mid %
  x^2 + y^2 - y \leq 0, x \geq 0\}$;
\item $\ds\iint_D \ln(1 + x^2 + y^2) dxdy$, unde $D$ este mărginit de curbele de ecuații:
  \[
    \begin{cases}
      x^2 + y^2 &= e^2 \\
      y &= x \sqrt{3} \\
      x &= y\sqrt{3} \\
      x &\geq 0
    \end{cases}
  \]
\end{enumerate}

\end{document}
